\section{Limits within the construction}
\label{sec:limits}

This section asks what the construction of
Sections~\ref{sec:stages} and~\ref{sec:chains} cannot do. It fixes the
class of schemes (Section~\ref{subsec:limits-class}), shows that
\(p>9/4\) even with compressed auxiliaries
(Section~\ref{subsec:limits-exponent}), bounds the copy counts and
\(\log(B/A)\) from below as \(p\to9/4\)
(Section~\ref{subsec:limits-resources}) and the number of target copies
for \(p<3\) (Section~\ref{subsec:limits-small}), and lists what the
class leaves out (Section~\ref{subsec:limits-scope}). None of these
statements is a lower bound for \(\omega\) or for tensor rank.

\subsection{Chain data and supplies}
\label{subsec:limits-class}

\begin{definition}[Chain data and supplies]\label{def:class}
\emph{Chain data} consist of the following.
\begin{enumerate}
\item[(a)] Integers \(J,W\ge1\), and pairs \((r_i,b_i)\) of positive
integers for \(1\le i\le J+1\).
\item[(b)] For each \(i\le J\), an arity \(d_i\ge2\) and a child length
\(c_i\ge1\) such that, for one ordering \((r,b)\) of the entries of
\((r_i,b_i)\), the stage \((r,b)\to(c_i,r)\) of arity \(d_i\) is legal
\eqref{eq:stage-conditions} and, after exchanging \(r_{i+1}\) and
\(b_{i+1}\) if necessary, \(r_{i+1}\le c_i\) and \(b_{i+1}\le r\).
\item[(c)] For each \(i\le J\), finitely many \emph{blocks} \(j\), each
with a power \(k_{ij}\ge1\) and a type \(\tau_{ij}\) of width
\(N_{ij}\ge1\) (Lemma~\ref{lem:stage-types}), such that
\(\sum_jN_{ij}k_{ij}=W\).
\item[(d)] A \emph{terminal} \(t\) with \(1\le t\le\max(r_{J+1},b_{J+1})\).
\end{enumerate}
Put \(s=r_1+b_1-1\). Let \(M_{ij}\) be the sector
count~\eqref{eq:stage-sector-count} of \(\tau_{ij}\) at arity \(d_i\),
\(E_{ij}=1+6k_{ij}(M_{ij}-1)\), \(U_i=d_i^{2W}\) and
\(V_i=\prod_jM_{ij}^{2k_{ij}}\), and put
\begin{equation}\label{eq:class-counts}
 \bar A=\prod_{i=1}^JU_i,\quad
 \bar B=t^{2W}\prod_{i=1}^JV_i,\quad
 r_*=s^{6W}\prod_{i,j}E_{ij}(2M_{ij}-1)^{6k_{ij}},\quad
 m_*=\prod_{i,j}M_{ij}^{6k_{ij}} .
\end{equation}
A \emph{supply} for chain data consists of positive integers
\(A=\mathfrak b_0\) and \(f_1,\ldots,f_J\) such that \(U_if_i\le\mathfrak b_{i-1}\)
for every \(i\), where \(\mathfrak b_i=f_iV_i\). Its \emph{source side} is
\(A\), and its \emph{target side} is \(B=t^{2W}\mathfrak b_J\).
\begin{itemize}
\item The \emph{source supply} has \(A=\bar A\) and
\(f_i=\prod_{l>i}U_l\prod_{l<i}V_l\).
\item The \emph{recycling supply} with source side \(A\) has
\(f_i=\lfloor\mathfrak b_{i-1}/U_i\rfloor\). It is a supply if every
\(f_i\ge1\).
\end{itemize}
By Lemma~\ref{lem:recycling}, chain data and a supply give a restriction
\(r_*\MM_A\restr m_*\MM_B\) over \(\Q\). If \(B>A\), it is a relative
scheme, which we call a \emph{chain scheme}.
\end{definition}

Block \(j\) of stage \(i\) applies the exact relation of
Corollary~\ref{cor:exact-stage} with \((N,k,\tau)=(N_{ij},k_{ij},\tau_{ij})\)
to its own \(N_{ij}k_{ij}\) of the \(W\) factors of
\(Q_{r_i,b_i}^{\otimes W}\). So stage \(i\) consumes an auxiliary of side
\(U_i\) and produces one of side \(V_i\), and \(\bar A\) is the total side
consumed. Before stage \(i\) the scheme holds one matrix multiplication
factor \(\MM_{\mathfrak b_{i-1}}\), the \emph{bank}; stage \(i\) cuts
\(\MM_{U_i}\MM_{f_i}\) out of it, turns \(\MM_{U_i}\) into
\(\MM_{V_i}\), and keeps \(\MM_{V_i}\MM_{f_i}\cong\MM_{\mathfrak b_i}\). Under
the source supply, \(U_if_i=\mathfrak b_{i-1}\) for every \(i\), and
\(B=\bar B\). Under the recycling supply each stage takes what it can
from the factors that earlier stages produced.

The class fixes the stage relation of Corollary~\ref{cor:exact-stage}
for every block (uncompressed auxiliaries, Lü's separator with \(2M-1\)
copies, separate exactification of each block), the evaluation of the
source with \(s^{6W}\) products, and the crop of the last pair to
\((1,t)\). The arities, children, crops, exchanges, blocks, types, widths,
powers, the terminal and the supply are free. The restrictions of
Proposition~\ref{prop:chain-counts} and Remark~\ref{rem:varying-powers}
are given by chain data with one block per stage and the source supply,
without crops or exchanges; so the chain schemes of
Section~\ref{sec:chains} are chain schemes in the sense of
Definition~\ref{def:class}.

\begin{lemma}[Supplies and forward recycling]\label{lem:recycling}
Let chain data and a supply be given, with source side \(A\) and target
side \(B\). Then \(r_*\MM_A\restr m_*\MM_B\) over \(\Q\). Moreover:
\begin{enumerate}
\item[(a)] \(B/A\le\bar B/\bar A\), with equality for the source supply;
\item[(b)] if \(B>A\), then this restriction and the one given by the
source supply are chain schemes, and the pooled exponent of the first is
at least that of the second;
\item[(c)] if \(A\ge4J\max_{1\le i\le J}U_i\prod_{l<i}(U_l/V_l)\), then
the recycling supply with source side \(A\) is a supply,
\(\mathfrak b_{i-1}\ge2JU_i\) for every \(i\), and \(B/A\ge\bar B/(2\bar A)\).
\end{enumerate}
\end{lemma}

The proof (Appendix~\ref{subsec:proof-recycling}) is an induction over
the stages. The bank only carries factors forward, so this is not a
catalyst, and every auxiliary keeps its full side. It works because
every stage is exactified before its output is used
(Section~\ref{subsec:prelim-degenerations}). The same induction can also
fund the blocks of one stage one after another, since they act on
disjoint factors of \(Q\). This keeps \(r_*\), \(m_*\) and
\(B/A\le\bar B/\bar A\), so the results of
Section~\ref{subsec:limits-resources} on \(r_*\), \(m_*\) and
\(\log(B/A)\) hold for it with the same proofs, as they use only
Lemma~\ref{lem:recycling}(a) and~(b).

\begin{remark}[Composition]\label{rem:composition}
The step of this induction composes any two relative schemes. Let
\(r^{(l)}\MM_{A^{(l)}}\restr m^{(l)}\MM_{B^{(l)}}\), \(l=1,2\), be
relative schemes with \(B^{(1)}\ge A^{(2)}\), and put
\(f=\lfloor B^{(1)}/A^{(2)}\rfloor\). Take \(r^{(2)}\) copies of the
first, restrict each \(\MM_{B^{(1)}}\) to its principal block
\(\MM_{A^{(2)}f}\cong\MM_{A^{(2)}}\MM_f\), and apply the second, tensored
with the identity on \(\MM_f\). This gives
\(r^{(1)}r^{(2)}\MM_{A^{(1)}}\restr m^{(1)}m^{(2)}\MM_{B^{(2)}f}\). Since
\(B^{(2)}f/A^{(1)}\le(B^{(1)}/A^{(1)})(B^{(2)}/A^{(2)})\) and
\(r^{(l)}/m^{(l)}>1\) (Section~\ref{subsec:prelim-schemes}), the
composite, if \(B^{(2)}f>A^{(1)}\), has pooled exponent at least the
smaller of the two.

The classical algorithm is the relative scheme \(n^3\MM_1\restr\MM_n\),
with source side \(1\) and pooled exponent \(3\). Let
\(r_*\MM_A\restr m_*\MM_B\) be a relative scheme, and let
\(n\ge2AB/(B-A)\). Put the classical algorithm first, and then apply
the scheme \(j\) times, each time to the factor produced by the previous
step. The composite has source side \(1\), copy counts \(n^3r_*^j\) and
\(m_*^j\), and target side \(\mathfrak b_j\), where \(\mathfrak b_0=n\) and
\(\mathfrak b_{l+1}=\lfloor\mathfrak b_l/A\rfloor B\). Every step can be applied, and
\(\log\mathfrak b_j\ge\log n+j\log(B/A)-2\)
(Appendix~\ref{subsec:proof-recycling}). So the pooled exponent of the
composite tends to that of the scheme as \(j\to\infty\), and with
Theorem~\ref{thm:main}, for every \(\delta>0\) some composite with
source side \(1\) has pooled exponent below \(9/4+\delta\).

So the source side alone does not measure the size of an algorithm. A
lower bound on \(A\), such as Theorem~\ref{thm:lower}(c) below,
concerns chain schemes with the source supply. The copy counts \(r_*\)
and \(m_*\), and hence \(r_*/m_*\), multiply under composition and
tensor products.
\end{remark}

\subsection{The exponent limit}
\label{subsec:limits-exponent}

\paragraph{Scalar arguments.}
Lemma~\ref{lem:profile-potential} is a scalar argument, the analogue
of~\citep[Remark~3.3]{FinCon26}: it bounds the gain
\(\log\frac{d+2}d\) of each stage by the drop of \(\log F\) along the
stage. A compressed auxiliary changes the gain of a stage, so the rest
of this subsection works with tensors.

\paragraph{Sector tensors.}
Let \(r,c,N\ge1\), and write \([r]=\{0,\ldots,r-1\}\). The \emph{middle
child} is
\[
 C^{\mathrm{mid}}(r,c)=\sum_{\alpha=0}^{r-1}\sum_{\gamma=0}^{c-1}
 x_\alpha\,y_{r-1-\alpha+\gamma}\,z_\gamma ,
\]
with legs of sizes \(r\), \(r+c-1\) and \(c\). Reversing the first
index turns it into \(C(r,c)^{(23)}\); we keep the first index
unreversed, because the bound below needs the first leg fixed. Let
\(X=(\mathbb F^r)^{\otimes N}\), with coordinates \(x_\alpha\) indexed by
words \(\alpha\in[r]^N\). An \emph{\((r,c;N)\) sector tensor} is a sum
\(T=\sum_{g=1}^MT_g\) of \(M\) sectors sharing the first leg \(X\)
(Section~\ref{subsec:prelim-sectors}) such that, for each \(g\),
invertible linear maps on \(Y_g\) and \(Z_g\), together with the
identity on \(X\), carry \(T_g\) to a product
\(L_{g,1}\otimes\cdots\otimes L_{g,N}\); here each \(L_{g,\ell}\) is
\(C(r,c)\) or \(C^{\mathrm{mid}}(r,c)\), and letter \(\ell\) of
\(\alpha\) is its first index. Let \(\nmid(g)\) be the number of factors
\(C^{\mathrm{mid}}(r,c)\). For tensors \(S\) and \(T\) with the same
first leg, we write \(S\restr_1T\) if linear substitutions on legs \(2\)
and \(3\), together with the identity on leg~\(1\), transform \(S\) into
\(T\).

Let \((r,b)\to(c,r)\) be a legal stage of arity \(d\), and let \(T\) be
the tensor of Lemma~\ref{lem:stage-types} for a type
\(\tau=(\nK,\nO,\nmid)\). Then \(T\) is an \((r,c;N)\) sector tensor with
\(M\) sectors given by~\eqref{eq:stage-sector-count}, and
\(\nmid(g)=\nmid\) for all \(g\): its sectors are
\(T_w=\bigotimes_\ell U'_{w_\ell}\), and by
Lemmas~\ref{lem:shifted-chain}(c) and~\ref{lem:reset}(b), renaming the
coordinates of legs \(2\) and \(3\) alone identifies \(U'_\chi\) with
\(C(r,c)\) for kernel and outer letters and with
\(C^{\mathrm{mid}}(r,c)\) for the middle letter. The substitutions of
Lemma~\ref{lem:shifted-chain}(c) and the zeroing of
Lemma~\ref{lem:stage-types} act on legs \(2\) and \(3\) only, so
\(C(r,b)^{\otimes N}B_X(d^N)\restr_1T\).

\paragraph{The Hankel lift.}
Let \(u,v\ge1\) with \(u+v=r+1\). For a tensor \(T\) with first leg
\(X\), write \(T[\alpha,\eta,\zeta]\) for its coefficient at
\(x_\alpha y_\eta z_\zeta\). Its \emph{Hankel lift} \(\mathcal L_u(T)\)
is the matrix with rows indexed by pairs \((\zeta,a)\), where \(\zeta\)
is a coordinate of leg \(3\) and \(a\in[u]^N\), columns indexed by pairs
\((\eta,e)\), where \(\eta\) is a coordinate of leg \(2\) and
\(e\in[v]^N\), and entries
\[
 \mathcal L_u(T)_{(\zeta,a),(\eta,e)}=T[a+e,\eta,\zeta],
\]
where \(a+e\in[r]^N\) is the word with letters \(a_\ell+e_\ell\).
Equivalently, \(\mathcal L_u(T)\) is an ordinary flattening of the
tensor obtained from \(T\) by pulling its first leg back along the
multiplication map \((\mathbb F^u\otimes\mathbb F^v)^{\otimes N}\to X\)
of \(C(u,v)^{\otimes N}\); this is a catalecticant (Hankel) flattening.
Three facts follow from the definition.
\begin{enumerate}
\item[(i)] If \(S\restr_1T\), with matrices \(\mathsf B\) and
\(\mathsf\Gamma\) of the substitutions on legs \(2\) and \(3\), then
\(\mathcal L_u(T)=(\mathsf\Gamma^{\top}\otimes I)\,\mathcal L_u(S)\,(\mathsf B\otimes I)\),
so \(\operatorname{rank}\mathcal L_u(T)\le\operatorname{rank}\mathcal L_u(S)\).
\item[(ii)] For tensors with first legs \((\mathbb F^r)^{\otimes N}\) and
\((\mathbb F^r)^{\otimes N'}\), \(\mathcal L_u(TT')\) is
\(\mathcal L_u(T)\otimes\mathcal L_u(T')\) up to permutations of rows and
columns.
\item[(iii)] \(\mathcal L_u\bigl(C(r,b)^{\otimes N}B_X(q)\bigr)\) is
\(\mathcal L_u(C(r,b)^{\otimes N})\otimes I_q\) up to permutations, since
the coefficient of \(B_X(q)\) at \(y_\eta z_\zeta\) is \([\eta=\zeta]\).
\end{enumerate}
Over every field, the rank of a Kronecker product is the product of the
ranks.

\begin{lemma}[Local ranks]\label{lem:hankel-ranks}
For \(N=1\) and over every field,
\(\operatorname{rank}\mathcal L_u(C(r,b))=b+v-1\) and
\(\operatorname{rank}\mathcal L_u(C^{\mathrm{mid}}(r,c))=c+u-1\).
\end{lemma}

The proof (Appendix~\ref{subsec:proof-hankel}) factors each lift through
the sum of two indices and exhibits an identity submatrix of full size.

\begin{theorem}[Hankel flattening bound]\label{thm:hankel}
Let \(T\) be an \((r,c;N)\) sector tensor with \(M\) sectors over a
field \(\mathbb F\), and let \(b,q\ge1\). If
\(C(r,b)^{\otimes N}B_X(q)\restr_1T\), then for all integers \(u,v\ge1\)
with \(u+v=r+1\),
\begin{equation}\label{eq:hankel-bound}
 q\,(b+v-1)^N\ \ge\ \sum_{g=1}^M(c+v-1)^{N-\nmid(g)}(c+u-1)^{\nmid(g)} .
\end{equation}
\end{theorem}

\begin{proof}
By (ii), (iii) and Lemma~\ref{lem:hankel-ranks}, the left side is the
rank of \(\mathcal L_u(C(r,b)^{\otimes N}B_X(q))\). Since
\(T_g\in X\otimes Y_g\otimes Z_g\), the matrix \(\mathcal L_u(T)\) is
block diagonal, with the block \(\mathcal L_u(T_g)\) on the rows
\(Z_g\times[u]^N\) and the columns \(Y_g\times[v]^N\), so its rank is
\(\sum_g\operatorname{rank}\mathcal L_u(T_g)\). By (i), applied to the
invertible maps and to their inverses, (ii) and
Lemma~\ref{lem:hankel-ranks}, the rank of \(\mathcal L_u(T_g)\) is the
\(g\)-th term on the right. Now apply (i) to
\(C(r,b)^{\otimes N}B_X(q)\restr_1T\).
\end{proof}

Write \(\varphi_r(y)=(r+2y-1)^2-1\) for real \(y\ge1\), and define the
potential \(\Psi(r,b)=r^2+4rb+b^2-2r-2b-\frac16\) for real \(r,b\ge1\).

\begin{lemma}[Gain per factor]\label{lem:hankel-gain}
In Theorem~\ref{thm:hankel}, if \(b\ge c\), then
\(\frac1N\log(M/q)\le\frac12\log\bigl(\varphi_r(b)/\varphi_r(c)\bigr)\).
\end{lemma}

\begin{lemma}[Potential]\label{lem:hankel-potential}
For real \(r,b\ge1\):
\begin{enumerate}
\item[(a)] \(\Psi\) is symmetric and increasing in each argument, and
\(\Psi(r,b)\ge\Psi(1,1)=\frac{11}6\);
\item[(b)] if \(b>c\ge1\), then
\(\frac12\log\bigl(\varphi_r(b)/\varphi_r(c)\bigr)<\frac23\log\bigl(\Psi(r,b)/\Psi(r,c)\bigr)\);
\item[(c)] \(\Psi(r,b)\le\frac{11}6(r+b-1)^2\) if \(r+b\ge3\), and
\(\Psi(r,b)\ge\frac{11}6b^{3/2}\).
\end{enumerate}
\end{lemma}

The proofs are in Appendix~\ref{subsec:proof-hankel}. For
Lemma~\ref{lem:hankel-gain}, one takes \(u=v\) if \(r\) is odd, and
combines the two nearly balanced choices of \((u,v)\) by the
Cauchy--Schwarz inequality if \(r\) is even. Part~(b) of
Lemma~\ref{lem:hankel-potential} says that
\(\frac23\log\Psi(r,y)-\frac12\log\varphi_r(y)\) increases in \(y\).

\begin{definition}[Compressed auxiliaries]\label{def:compressed}
Let \((r,b)\to(c,r)\) be a legal stage of arity \(d\), let \(N\ge1\),
let \(\tau\) be a type of width \(N\), and let \(T\) be the tensor of
Lemma~\ref{lem:stage-types}. A \emph{compressed auxiliary of rank \(q\)}
for these data is a restriction \(C(r,b)^{\otimes N}B_X(q)\restr_1T\),
where \(q\ge1\). \emph{Compressed chain data} are chain data together
with, for every block \((i,j)\), a compressed auxiliary of some rank
\(q_{ij}\), used in all six orientations, for the stage of
Definition~\ref{def:class}(b) and \((N_{ij},\tau_{ij})\); in their
counts~\eqref{eq:class-counts}, \(U_i=\prod_jq_{ij}^{2k_{ij}}\).
\end{definition}

The stage itself has \(q=d^N\). Using a compressed auxiliary in every
orientation of the proof of Proposition~\ref{prop:stage} and then
Corollary~\ref{cor:exact-stage} gives, since
\(\bigotimes_\pi B_X(q)^\pi\cong\MM_{q^2}\),
\[
 E\,(2M-1)^{6k}\,Q_{r,b}^{\otimes Nk}\,\MM_{q^{2k}}\restr
 M^{6k}\,Q_{c,r}^{\otimes Nk}\,\MM_{M^{2k}}
\]
over every field of characteristic zero that contains the coefficients
of the compressed auxiliary. So Lemma~\ref{lem:recycling} and its proof
hold for compressed chain data, with the new \(U_i\).

\begin{corollary}[Compression limit]\label{cor:compression-limit}
Let compressed chain data and a supply give a relative scheme
\(r_*\MM_A\restr m_*\MM_B\). Then its pooled exponent satisfies
\(p>9/4\).
\end{corollary}

\begin{proof}
By Lemma~\ref{lem:recycling}(b), we may assume the source supply, so
\(A=\bar A\) and \(B=\bar B\). Put \(w_{ij}=N_{ij}k_{ij}/W\), so that
\(\sum_jw_{ij}=1\). Then
\[
 \frac{\log(B/A)}{2W}=\log t+\sum_{i,j}w_{ij}h_{ij},\qquad
 h_{ij}=\frac1{N_{ij}}\log\frac{M_{ij}}{q_{ij}} .
\]
Fix \(i\), and let \((r,b)\to(c_i,r)\) be the stage of
Definition~\ref{def:class}(b); it has \(c_i<b\). By
Theorem~\ref{thm:hankel} and Lemmas~\ref{lem:hankel-gain}
and~\ref{lem:hankel-potential}(b), applied to each block, and by the
symmetry of \(\Psi\),
\(\sum_jw_{ij}h_{ij}<\frac23\log\bigl(\Psi(r_i,b_i)/\Psi(c_i,r)\bigr)\).
By Lemma~\ref{lem:hankel-potential}(a),
\(\Psi(r_{i+1},b_{i+1})\le\Psi(c_i,r)\) and
\(\Psi(1,t)\le\Psi(r_{J+1},b_{J+1})\). So the sum telescopes, and
Lemma~\ref{lem:hankel-potential}(c), with \(r_1+b_1\ge3\), gives
\[
 \frac{\log(B/A)}{2W}<\log t+\frac23\log\frac{\Psi(r_1,b_1)}{\Psi(1,t)}
 \le\log t+\frac23\log\frac{s^2}{t^{3/2}}=\frac43\log s .
\]
By the proof of Lemma~\ref{lem:recycling}(b),
\(\log(r_*/m_*)\ge6W\log s>0\). Hence
\(p=\log(r_*/m_*)/\log(B/A)>6/(8/3)=9/4\).
\end{proof}

For \(q_{ij}=d_i^{N_{ij}}\), the corollary gives the conclusion
\(p>9/4\) of Lemma~\ref{lem:profile-potential}(c) for every chain scheme,
with any supply, but not the intermediate bound
\(p>3\log s/(\log t+G)\). One rank in all six orientations is no loss of
generality for the barrier: if orientation \(\pi\) has a compressed
auxiliary of rank \(q_\pi\), then \(B_X(\max_\pi q_\pi)\restr B_X(q_\pi)\)
(set the extra coordinates to zero) gives one of rank \(\max_\pi q_\pi\)
in every orientation, as the square auxiliary \(\MM_{q^2}\) requires,
and~\eqref{eq:hankel-bound} holds for each \(q_\pi\) separately.

\begin{remark}[Ordinary sectors only]\label{rem:ordinary-only}
Middle sectors are what make exponents below \(3\) possible. For
\(u=1\), \(\mathcal L_1(T)\) is the flattening of \(T\) with rows indexed
by leg~\(3\), whose rank does not increase under any restriction, and if
no sector has a middle factor, then~\eqref{eq:hankel-bound} reads
\(M(r+c-1)^N\le q(r+b-1)^N\). So let compressed chain data have only
types with \(\nmid=0\), and put \(s_i=r_i+b_i-1\), so that \(s_1=s\). In
the notation of the proof above, every block of stage \(i\) has
\(h_{ij}\le\log(s_i/(r+c_i-1))\le\log(s_i/s_{i+1})\). These bounds
telescope, and
\(t\le s_{J+1}\) gives \(\log(B/A)\le2W\log s\). Hence \(p\ge3\) for
every supply.
\end{remark}

\begin{remark}[The first leg, and sharpness]\label{rem:hankel-remarks}
For \(u=1\) and \(u=r\), the lift is a flattening of \(T\) along leg
\(3\) or leg \(2\). The lifts with \(2\le u\le r-1\) use the Hankel
structure of the shared first leg, so they need \(\restr_1\). Crops of
the first leg to a prefix are harmless: they turn \(C(r,b)\) into
\(C(r',b)\) with unused coordinates, and \(\Psi\) increases in \(r\).
General first-leg maps are not covered. Theorem~\ref{thm:hankel} also
holds if \(T\) is the constant term of a degeneration with the identity
on leg~\(1\): the minors of size \(q(b+v-1)^N+1\) of the lift of the
transformed tensor vanish as polynomials in the parameter, hence at
\(0\). So removing cross terms by powers of a parameter does not escape
the bound. In the homogeneous limit, Lemma~\ref{lem:hankel-gain} is
tight on the efficient ray \(r=\lambda b\), \(c=\lambda^2b\) with
\((d+2)\lambda^2+\lambda-d=0\), where \((r+2b)/(r+2c)=(d+2)/d\), the gain
per factor of an uncompressed stage. The potential step of
Lemma~\ref{lem:hankel-potential}(b) is tight only to first order as
\(\lambda\to1\). For compressed auxiliaries this suggests, but we do not
prove, that \(9/4\) can be approached only with growing arity; for chain
schemes, which have uncompressed auxiliaries, Theorem~\ref{thm:lower}
proves it.
\end{remark}

\subsection{Resource bounds}
\label{subsec:limits-resources}

\begin{theorem}[Resource lower bound within the construction]\label{thm:lower}
Let \(0<\eps<1/9216\), and let a chain scheme \(r_*\MM_A\restr m_*\MM_B\)
have pooled exponent at most \(9/4+\eps\). Then, with \(\bar A\) as
in~\eqref{eq:class-counts}:
\begin{enumerate}
\item[(a)] \(\log r_*>\log m_*\ge3\log\bar A>3.9\cdot10^{-6}\,\eps^{-9/2}\log^3(1/\eps)\);
\item[(b)] \(\log(B/A)>9.8\cdot10^{-5}\,\eps^{-4}\log^2(1/\eps)\);
\item[(c)] under the source supply,
\(\log A=\log\bar A>1.3\cdot10^{-6}\,\eps^{-9/2}\log^3(1/\eps)\).
\end{enumerate}
Moreover, some block \((i,j)\) has \(d_i>(96\eps)^{-1/2}\),
\(N_{ij}>9d_i\log(1/\eps)/(16\eps)\) and
\(k_{ij}>d_i\log d_i/(4\eps)\), and
\(W>3\log^2(1/\eps)/(8192\,\eps^3)\).
\end{theorem}

The proof is in Appendix~\ref{subsec:lower-main-proof}. In~(a),
\(\bar A\) is the auxiliary side consumed by the stages; the supplied
source side \(A\) may be smaller (Corollary~\ref{cor:recycled-main}) or
larger, and (a) bounds \(r_*\) and \(m_*\) in every case.

The bounds \(3.9\cdot10^{-6}\,\eps^{-9/2}\log^3(1/\eps)\) on
\(\log r_*\) and \(\log m_*\) in~(a), and the bound~(b), also hold for
every relative scheme built from chain schemes and classical algorithms
by composition (Remark~\ref{rem:composition}) and tensor products.
Indeed, under both operations \(m_*\) and \(r_*/m_*\) are the products
of those of the parts, \(B/A\) is at most the product, and every part
has \(r_*/m_*>1\). So if the result has
\(\log(r_*/m_*)\le(9/4+\eps)\log(B/A)\), the mediant inequality gives a
part with \(\log(r_*/m_*)\le(9/4+\eps)\log(B/A)\), where
\(9/4+\eps<3\). This part is a chain scheme, and its \(m_*\) and
\(r_*/m_*\), hence also its \(r_*\), bound those of the result from
below. For~(b), the result has
\(\log(B/A)\ge\log(r_*/m_*)/(9/4+\eps)\), and the proof of~(b) bounds
this quotient from below for that part.

\begin{corollary}[Recycling the main construction]\label{cor:recycled-main}
Let \(0<\eps\le1\) be rational, and take the chain, type width, power and
types of Theorem~\ref{thm:main}. The recycling supply with source side
\(A'=4JD^{2Nk}\) is a supply. It gives a relative scheme
\(r_*\MM_{A'}\restr m_*\MM_{B'}\) over \(\Q\) with the same \(r_*\) and
\(m_*\) as Theorem~\ref{thm:main}, whose pooled exponent \(p'\) satisfies
\(p\le p'<p+10^{-9}\eps\), and so \(9/4<p'<9/4+\eps/2\). Moreover,
\[
 \log A'<5.32\cdot10^6\,\eps^{-3}\log^3\frac{4096}\eps,\qquad
 \log B'<1.5\cdot10^9\,\eps^{-4}\log^2\frac{4096}\eps .
\]
\end{corollary}

The proof is in Appendix~\ref{subsec:lower-recycled-proof}. For
example, on the chain of Theorem~\ref{thm:main} at \(\eps=1\), the source
side \(A=\prod_id_i^{2Nk}\) has \(3{,}250{,}941{,}263\) decimal digits
and the recycled source side \(A'=4JD^{2Nk}\) has \(12{,}783{,}943\)
(supplementary script \path|shifted_chains_check_recycling.py|,
Part~5).

\begin{remark}[Orders]\label{rem:limits-orders}
Within the construction, as \(\eps\to0\):
\begin{itemize}
\item \emph{Copy counts.} For every chain scheme, and for the composites
above, \(p\le9/4+\eps\) with \(\eps<1/9216\) forces \(\log m_*\) and
\(\log r_*\) above \(3.9\cdot10^{-6}\eps^{-9/2}\log^3(1/\eps)\). For
rational \(\eps\le\frac12\), Corollary~\ref{cor:main-eps} gives
\(p<9/4+\eps\) with \(m_*\le A^6\) and \(r_*<A^{35/4}\), so
\(\log m_*<8.4\cdot10^8\eps^{-9/2}\log^3(2048/\eps)\) and
\(\log r_*<1.23\cdot10^9\eps^{-9/2}\log^3(2048/\eps)\). So the least
values of \(\log m_*\) and \(\log r_*\) have order
\(\eps^{-9/2}\log^3(1/\eps)\).
\item \emph{Ratio.} The least \(\log(B/A)\) has order
\(\eps^{-4}\log^2(1/\eps)\), by Theorem~\ref{thm:lower}(b) and by the
scheme of Corollary~\ref{cor:main-eps}, for which the proof of
Corollary~\ref{cor:recycled-main}, with \(2\eps\) in place of \(\eps\), gives
\(\log(B/A)<1.43\cdot10^9(2\eps)^{-4}\log^2(2048/\eps)
<9\cdot10^7\eps^{-4}\log^2(2048/\eps)\). Recycling keeps this order,
since \(B/(2A)\le B'/A'\le B/A\) in the proof of
Corollary~\ref{cor:recycled-main}.
\item \emph{Source side.} Under the source supply, the least \(\log A\)
has order \(\eps^{-9/2}\log^3(1/\eps)\), by Theorem~\ref{thm:lower}(c)
and Corollary~\ref{cor:main-eps}. Otherwise the counts do not determine
it: Corollary~\ref{cor:recycled-main} gives
\(\log A'=O(\eps^{-3}\log^3(1/\eps))\), and composition with the
classical algorithm gives source side \(1\). Part~(c) needs its
hypothesis: for rational \(\eps\le2^{-26}\), the scheme of
Corollary~\ref{cor:recycled-main} has pooled exponent below
\(9/4+\eps/2\) and a source side below the bound of~(c) with
\(\eps/2\) in place of \(\eps\), since the ratio of the two bounds,
\(\frac{5.32\cdot10^{12}}{1.3}\,2^{-9/2}\eps^{3/2}\bigl(\log(4096/\eps)/\log(2/\eps)\bigr)^3\),
increases with \(\eps\) and is below \(0.92\) at \(\eps=2^{-26}\).
\item \emph{Seeds.} The seed of Theorem~\ref{thm:seed} depends on the
scheme through \(\Xi=r_*(1+r_*m_*)n_*^3\), and
\(r_*^2m_*\le\Xi\le54r_*^2m_*A^3\), where
\(r_*^2m_*\ge m_*^3=B^9\ge(2A)^9\) (Lemma~\ref{lem:resources}). So the
seed built in the same way from the scheme of
Corollary~\ref{cor:recycled-main} has \(\log n_0\) smaller by a factor
less than \(\frac43+\frac{\log55}{9\log A}<1.34\): the copy counts, not
the source side, set the size of the seed.
\end{itemize}
\end{remark}

\subsection{Small batches}
\label{subsec:limits-small}

Unlike the source side, the number \(m_*\) of target copies does not
depend on the supply, and it multiplies under composition. It is the
batch size of the pooled recursion~\citep[Lemma~2.4]{FinCon26}, and the
maps of a relative scheme have at least \(m_*B^2\) nonzero coefficients
on each leg (see the proof of Proposition~\ref{prop:small-sources}).

\begin{proposition}[Small batches; computer-assisted]\label{prop:small-sources}
Let \(M_0=\frac{23!}{6!\,11!\,6!}\,2^{11}=2\,558\,607\,507\,456\); then
\(2^{494}<M_0^{12}<2^{495}\).
\begin{enumerate}
\item[(a)] Every chain scheme with the source supply and \(p<3\) has
\(m_*\ge M_0^{12}\).
\item[(b)] So does every relative scheme with \(p<3\) obtained from such
chain schemes and the classical restrictions \(n^3\MM_1\restr\MM_n\) by
repeated use of these operations: (i) replacing \(r\MM_A\restr m\MM_B\)
by some \(r\MM_{A'}\restr m\MM_{B'}\) with \(B'/A'\le B/A\), as other
supplies do (Lemma~\ref{lem:recycling}); (ii) composition as in
Remark~\ref{rem:composition}, applied to restrictions
\(r^{(l)}\MM_{A^{(l)}}\restr m^{(l)}\MM_{B^{(l)}}\), \(l=1,2\), with
\(B^{(1)}\ge A^{(2)}\); and (iii) tensor products. The maps of such a
scheme have more than \(2^{496}\) nonzero coefficients on each leg.
\item[(c)] The bound in (a) is attained: the one-stage chain
\((361,182)\to(1,361)\) of arity \(2\) with \(N=23\), \(k=2\) and type
\((6,11,6)\) has \(m_*=M_0^{12}\) and \(p<2.9999985\).
\end{enumerate}
\end{proposition}

The proof is in Appendix~\ref{subsec:lower-small-proof}; it checks two
finite lists of cases with the supplementary script
\path|shifted_chains_check_small_sources.py|. The bound uses the charges
in~\eqref{eq:class-counts} and uncompressed auxiliaries. It says nothing
about schemes outside the construction, such as Strassen's, with
\(m_*=1\), and it concerns only the operations listed in~(b), since
keeping some of the target copies is a restriction too.

\subsection{Scope}
\label{subsec:limits-scope}

The results above cover the class of Definition~\ref{def:class},
compressed auxiliaries for the exponent, and supplies and composition
for \(r_*\), \(m_*\) and \(\log(B/A)\). For compressed chain data we
have no resource bound. The bounds concern counts only, so they do not
depend on the field.

The proofs extend to two variants of the class.
\begin{itemize}
\item \emph{Other separators and exactifiers.} The proofs of
Corollary~\ref{cor:compression-limit} and Theorem~\ref{thm:lower} use
the separator and the exactifier of a block \((i,j)\) only through their
numbers of copies. Corollary~\ref{cor:compression-limit} needs only that
the separator has at least \(M_{ij}\) copies, and
Theorem~\ref{thm:lower} needs this for the exactifier as well. For
the tensor \(T\) of Lemma~\ref{lem:stage-types}, with \(M\) sectors,
every separator with output~\eqref{eq:prelim-separator} has at least
\(M\) copies. Indeed, maps on legs~\(2\) and~\(3\) identify each sector
with a product of \(N\) factors \(C(r,c)\) or \(C^{\mathrm{mid}}(r,c)\)
(Section~\ref{subsec:limits-exponent}), each of rank \(r\) along
leg~\(1\). So the flattening of \(\bigoplus_gT_gB_X(M)\) along leg~\(1\)
has rank \(Mr^N\). That of \(cT\) has rank at most \(cr^N\), and
degenerations do not raise flattening ranks, so
\(cT\degen\bigoplus_gT_gB_X(M)\) needs \(c\ge M\).
Proposition~\ref{prop:small-sources} uses the exact counts
\(2M_{ij}-1\) and \(E_{ij}\).
\item \emph{Joint exactification of the blocks of one stage} keeps
Theorem~\ref{thm:lower} except for the bound on \(k_{ij}\)
(Remark~\ref{rem:joint-exactification}).
\end{itemize}

The class charges \(s^{6W}\) products for the source. Since the
flattening of \(Q_{r,b}\) along leg~\(1\) has rank \((rb(r+b-1))^2\),
every evaluation of \(Q_{r_1,b_1}^{\otimes W}\) needs at least
\((r_1b_1s)^{2W}\) products, and with such a source the proof of
Corollary~\ref{cor:compression-limit} gives only
\(p>\frac34\bigl(1+\log(r_1b_1)/\log s\bigr)\), which may be below
\(9/4\) for a finite chain.

Outside the class are:
\begin{enumerate}
\item[(i)] exactification of several stages together, and exactifiers
with fewer than \(M_{ij}\) copies;
\item[(ii)] auxiliaries shared between orientations or stages, or cut
from the \(Q\)-factors;
\item[(iii)] blocks of one stage with different children, factors that
skip a stage, and states that split and merge;
\item[(iv)] other local maps, such as the determinant splits
of~\citep[Section~4]{FinCon26}, children longer than their parents (for
which Lemma~\ref{lem:hankel-potential}(b) fails), stages of arity~\(1\),
and different types in different orientations;
\item[(v)] other terminal restrictions, and combinations with other
schemes beyond Remark~\ref{rem:composition}.
\end{enumerate}
